\section{Curvas en el plano y en el espacio}
\begin{ejercicio}
    Sea $\alpha:I\to \mathbb{R}^2$ una curva regular plana y $a\in \mathbb{R}^2\setminus \alpha(I)$. Si existe $t_0\in I$ de forma que $|\alpha(t)-a|\geq |\alpha(t_0)-a|$ para cada $t\in I$, prueba que la recta que une el punto $a$ con $\alpha(t_0)$ es la recta normal de $\alpha$ en $t_0$. Lo mismo también es cierto si le damos la vuelta a la desigualdad.\\

    \noindent
    La recta normal a $\alpha$ en $t_0$ es la recta:
    \begin{equation*}
        R = \alpha(t_0) + \langle J\alpha'(t_0) \rangle 
    \end{equation*}
    Ahora, si definimos $f:I\to \mathbb{R}$ dada por:
    \begin{equation*}
        f(t) = |\alpha(t)-a|^2, \qquad t\in I
    \end{equation*}
    La condición del enunciado nos dice que $f(t)\geq f(t_0)\quad \forall t\in I$, por lo que $f$ alcanza un mínimo en $t_0\in I$. Vemos que $f$ es diferenciable, con derivada:
    \begin{equation*}
        f'(t) = 2\langle \alpha(t)-a,\alpha'(t) \rangle  \qquad \forall t\in I
    \end{equation*}
    Como $f$ alcanza un mínimo en $t_0$ tenemos que:
    \begin{equation*}
        0 = f'(t_0) = 2\langle \alpha(t_0)-a,\alpha'(t_0) \rangle 
    \end{equation*}
    Es decir, $\alpha'(t_0)\perp \alpha(t_0)-a$, por lo que ha de existir $\lm\in \mathbb{R}$ de forma que:
    \begin{equation*}
        \alpha(t_0) - a = \lm J\alpha'(t_0)
    \end{equation*}

    de donde:
    \begin{equation*}
        a = \alpha(t_0) - \lm J \alpha'(t_0)
    \end{equation*}
    Por lo que $a\in R$. Si le damos la vuelta a la igualdad tendríamos entonces que $t_0\in I$ sería un máximo de $f$, y por el mismo razonamiento llegamos también a la misma tesis.
\end{ejercicio}

\begin{ejercicio}
    Sea $\alpha:I\to \mathbb{R}^2$ una curva regular plana y sea $R$ una recta de $\mathbb{R}^2$. Si uno puede encontrar un instante $t_0\in I$ de forma que la distancia de $\alpha(t)$ a $R$ sea mayor o igual que la distancia de $\alpha(t_0)$ a $R$ para todo $t\in I$ y de forma que $\alpha(t_0)\notin R$, entonces prueba que la recta tangente a $\alpha$ en $t_0$ es paralela a $R$.\\

    % \noindent
    % $R$ será de la forma:
    % \begin{equation*}
    %     R  = a + \langle v \rangle 
    % \end{equation*}
    % para ciertos $a\in \mathbb{R}^2,v\in \mathbb{R}^2\setminus \{0\}$. La recta tangente a $\alpha$ en $t_0$ es la recta:
    % \begin{equation*}
    %     S = \alpha(t_0) + \langle \alpha'(t_0) \rangle 
    % \end{equation*}
\end{ejercicio}
